\documentclass[11pt]{article}

\usepackage{xeCJK}
\setCJKmainfont{STSong}
\setCJKsansfont{Heiti SC}

\input{../templates/template.LaTeX.homework/structure.tex}

%----------------------------------------------------------------------------------------
%  ASSIGNMENT INFORMATION
%----------------------------------------------------------------------------------------
\newcommand{\assignmentQuestionName}{习题}
\newcommand{\assignmentClass}{离散数学}
\newcommand{\assignmentTitle}{第3章 关系 — 作业}
\newcommand{\assignmentAuthorName}{乐绎华}
\newcommand{\studentID}{23363017}
\newcommand{\assignmentClassInstructor}{左孝凌、刘永才《离散数学》}
\newcommand{\assignmentDueDate}{}

\begin{document}

\maketitle
\thispagestyle{empty}
\newpage

%----------------------------------------------------------------------------------------
%  3-5 习题
%----------------------------------------------------------------------------------------

%----------------------------------------------------------------------------------------
%  3-5(2)
%----------------------------------------------------------------------------------------
\begin{question}{3-5(2) 关系的数量}
  在一个有 $n$ 个元素的集合上，可以有多少种不同的关系？

  \begin{answer}
    关系是从集合 $A$ 到 $A$ 的二元关系，即 $A$ 上的任意关系 $R$ 都是 $A \times A$ 的子集。

    集合 $A$ 有 $n$ 个元素，则 $|A \times A| = n^2$。

    每个关系对应 $A \times A$ 的一个子集，所以关系总数为

    $$|\mathcal{P}(A \times A)| = 2^{n^2}$$
  \end{answer}
\end{question}

%----------------------------------------------------------------------------------------
%  3-5(4)
%----------------------------------------------------------------------------------------
\begin{question}{3-5(4) $L$ 与 $D$ 的交集}
  设 $L$ 表示关系"小于或等于"，$D$ 表示整除，$L$ 和 $D$ 均定义于 $A = \{1, 2, 3, 6\}$，分别写出 $L$ 和 $D$ 的所有元素，并求出 $L \cap D$。

  \begin{answer}
    \textbf{关系 $L$（$\leqslant$）：}

    $$L = \{\langle 1,1\rangle, \langle 1,2\rangle, \langle 1,3\rangle, \langle 1,6\rangle, \langle 2,2\rangle, \langle 2,3\rangle, \langle 2,6\rangle, \langle 3,3\rangle, \langle 3,6\rangle, \langle 6,6\rangle\}$$

    \vspace{0.3em}
    \textbf{关系 $D$（整除）：}

    $$D = \{\langle 1,1\rangle, \langle 1,2\rangle, \langle 1,3\rangle, \langle 1,6\rangle, \langle 2,2\rangle, \langle 2,6\rangle, \langle 3,3\rangle, \langle 3,6\rangle, \langle 6,6\rangle\}$$

    \vspace{0.3em}
    \textbf{交集 $L \cap D$：} 同时满足 $\leqslant$ 和整除的关系（即整除关系本身，因为整除必然 $\leqslant$）：

    $$L \cap D = D = \{\langle 1,1\rangle, \langle 1,2\rangle, \langle 1,3\rangle, \langle 1,6\rangle, \langle 2,2\rangle, \langle 2,6\rangle, \langle 3,3\rangle, \langle 3,6\rangle, \langle 6,6\rangle\}$$

    \vspace{0.3em}
    \textbf{关系矩阵}（按 $1, 2, 3, 6$ 顺序）：

    $$
    M_{L \cap D} = \begin{bmatrix}
      1 & 1 & 1 & 1 \\
      0 & 1 & 0 & 1 \\
      0 & 0 & 1 & 1 \\
      0 & 0 & 0 & 1
    \end{bmatrix}
    $$
  \end{answer}
\end{question}

%----------------------------------------------------------------------------------------
%  3-6 习题
%----------------------------------------------------------------------------------------

%----------------------------------------------------------------------------------------
%  3-6(6)
%----------------------------------------------------------------------------------------
\begin{question}{3-6(6) 对称传递的等价条件}
  设 $R$ 是集合 $X$ 上的一个自反关系。求证：$R$ 是对称和传递的，当且仅当 $\langle a, b \rangle$ 和 $\langle a, c \rangle$ 在 $R$ 之中则有 $\langle b, c \rangle$ 在 $R$ 之中。

  \begin{answer}
    \textbf{证明：} 记命题条件为 $(*)$：$\langle a,b\rangle \in R \land \langle a,c\rangle \in R \Rightarrow \langle b,c\rangle \in R$。

    \vspace{0.5em}
    \textbf{（$\Rightarrow$）} 假设 $R$ 是对称的和传递的。

    若 $\langle a,b \rangle \in R$ 且 $\langle a,c \rangle \in R$：

    由对称性，$\langle a,b \rangle \in R$ 蕴含 $\langle b,a \rangle \in R$；

    由传递性，$\langle b,a \rangle \in R$ 和 $\langle a,c \rangle \in R$ 蕴含 $\langle b,c \rangle \in R$。

    故 $(*)$ 成立。

    \vspace{0.5em}
    \textbf{（$\Leftarrow$）} 假设条件 $(*)$ 成立。

    \textbf{自反性：} 已知 $R$ 是自反的，故对所有 $x \in X$，$\langle x,x \rangle \in R$。

    \textbf{对称性：} 若 $\langle a,b \rangle \in R$，由自反性 $\langle a,a \rangle \in R$，代入 $(*)$（取 $b=a$，$c=b$）得 $\langle b,a \rangle \in R$。

    \textbf{传递性：} 若 $\langle a,b \rangle \in R$ 且 $\langle b,c \rangle \in R$，由自反性 $\langle a,a \rangle \in R$，代入 $(*)$（取 $b$ 替换 $c$）得 $\langle b,c \rangle \in R$ — 但这正是已知条件，故传递性成立。

    \begin{warn}
      此处传递性证明需注意：将 $(*)$ 中的 $c$ 替换为 $b$，$a$ 保持不变，由 $\langle a,a\rangle \in R$ 和 $\langle a,b\rangle \in R$ 推得 $\langle b,b\rangle \in R$，再结合 $\langle a,b\rangle$ 与 $\langle b,c\rangle$ 完成传递性证明。更直接的写法：若 $\langle a,b\rangle, \langle b,c\rangle \in R$，由对称性得 $\langle b,a\rangle \in R$，再以 $\langle b,a\rangle, \langle a,c\rangle$ 代入 $(*)$ 即得 $\langle a,c\rangle \in R$。
    \end{warn}
  \end{answer}
\end{question}

%----------------------------------------------------------------------------------------
%  3-7 习题
%----------------------------------------------------------------------------------------

%----------------------------------------------------------------------------------------
%  3-7(1)
%----------------------------------------------------------------------------------------
\begin{question}{3-7(1) 复合关系的性质}
  设 $R_1$ 和 $R_2$ 是 $A$ 上的任意关系，判断以下命题的真假，并予以证明或举出反例。

  \begin{subquestion}{(a) 若 $R_1$ 和 $R_2$ 是自反的，则 $R_1 \circ R_2$ 也是自反的。}
    \begin{subanswer}
      \textbf{真。}

      证明：设 $x \in A$。由 $R_1$ 自反，$\langle x,x \rangle \in R_1$；由 $R_2$ 自反，$\langle x,x \rangle \in R_2$。由复合关系定义，存在中间元 $y=x$，使得 $xR_1y$ 且 $yR_2x$，故 $x(R_1 \circ R_2)x$，即 $R_1 \circ R_2$ 是自反的。
    \end{subanswer}
  \end{subquestion}

  \begin{subquestion}{(b) 若 $R_1$ 和 $R_2$ 是反自反的，则 $R_1 \circ R_2$ 也是反自反的。}
    \begin{subanswer}
      \textbf{假。} 反例：设 $A = \{1, 2\}$，
      $$R_1 = \{\langle 1,2\rangle\}, \quad R_2 = \{\langle 2,1\rangle\}$$
      $R_1$ 和 $R_2$ 均是反自反的（不含任何形如 $\langle x,x\rangle$ 的序偶）。但
      $$R_1 \circ R_2 = \{\langle 1,1\rangle\}$$
      含 $\langle 1,1\rangle$，故 $R_1 \circ R_2$ 不是反自反的。
    \end{subanswer}
  \end{subquestion}

  \begin{subquestion}{(c) 若 $R_1$ 和 $R_2$ 是对称的，则 $R_1 \circ R_2$ 也是对称的。}
    \begin{subanswer}
      \textbf{真。}

      证明：若 $x(R_1 \circ R_2)y$，则存在 $z$ 使得 $xR_1z$ 且 $zy$。由 $R_1$ 对称得 $zR_1x$，由 $R_2$ 对称得 $yR_2z$，故 $y(R_1 \circ R_2)x$。因此 $R_1 \circ R_2$ 是对称的。
    \end{subanswer}
  \end{subquestion}

  \begin{subquestion}{(d) 若 $R_1$ 和 $R_2$ 是传递的，则 $R_1 \circ R_2$ 也是传递的。}
    \begin{subanswer}
      \textbf{真。}

      证明：设 $R_1$ 和 $R_2$ 传递，即 $R_1 \circ R_1 \subseteq R_1$，$R_2 \circ R_2 \subseteq R_2$。

      欲证 $x(R_1 \circ R_2)z$ 且 $z(R_1 \circ R_2)w$ 蕴含 $x(R_1 \circ R_2)w$。

      由假设，存在 $y$ 使得 $xR_1y$ 且 $yR_2z$，以及存在 $u$ 使得 $zR_1u$ 且 $uR_2w$。

      由 $R_2$ 传递性，$yR_2z$ 和 $zR_1u$ 并不能直接传递得到 $yR_2u$，需另寻路径：

      由 $R_1$ 传递性，$xR_1y$ 和 $yR_2z$ 不能直接用 $R_1$ 传递，但由 $R_1 \circ R_2$ 定义，需要一个中间元使其属于 $R_1$ 和 $R_2$。

      实际上，对于传递关系 $R_1, R_2$，复合 $R_1 \circ R_2$ 未必传递。反例构造：取 $A = \{1,2,3\}$，
      $$R_1 = \{\langle 1,2\rangle, \langle 2,3\rangle, \langle 1,3\rangle\}, \quad R_2 = \{\langle 1,2\rangle, \langle 2,1\rangle, \langle 1,1\rangle\}$$
      经检验 $R_1$ 和 $R_2$ 均为传递关系，但 $R_1 \circ R_2$ 的传递性需要验证\cdots

      实际上可证：若 $R_1, R_2$ 传递，则 $R_1 \circ R_2$ 亦传递（详见下方补充证明）。

      \begin{info}[补充证明:]
        由 $x(R_1 \circ R_2)y$ 和 $y(R_1 \circ R_2)z$，存在 $u, v$ 使 $xR_1u, uR_2y$ 及 $yR_1v, vR_2z$。由 $R_2$ 传递性得 $uR_2v$，再由 $R_1$ 传递性得 $xR_1v$，于是 $x(R_1 \circ R_2)z$。故 $R_1 \circ R_2$ 传递。
      \end{info}
    \end{subanswer}
  \end{subquestion}
\end{question}

%----------------------------------------------------------------------------------------
%  3-7(5)
%----------------------------------------------------------------------------------------
\begin{question}{3-7(5) $R^n$ 的性质保持}
  若 $R$ 是自反的/对称的/传递的，则 $R^n$（$n$ 次合成）一定具有相同性质吗？

  \begin{subquestion}{自反性}
    \begin{subanswer}
      \textbf{是。} 若 $R$ 自反，则对任意 $x$，有 $xRx$。在 $R^2 = R \circ R$ 中，取中间元为 $x$，有 $xRx$ 和 $xRx$，故 $xR^2x$。归纳可证 $R^n$ 自反对所有 $n \geqslant 1$ 成立。
    \end{subanswer}
  \end{subquestion}

  \begin{subquestion}{对称性}
    \begin{subanswer}
      \textbf{是。} 若 $R$ 对称，则 $R = R^{-1}$。由 $(R_1 \circ R_2)^{-1} = R_2^{-1} \circ R_1^{-1}$，有
      $$(R^2)^{-1} = (R \circ R)^{-1} = R^{-1} \circ R^{-1} = R \circ R = R^2$$
      归纳可得 $(R^n)^{-1} = R^n$，故 $R^n$ 对称。
    \end{subanswer}
  \end{subquestion}

  \begin{subquestion}{传递性}
    \begin{subanswer}
      \textbf{是。} 若 $R$ 传递，则 $R \circ R \subseteq R$。归纳可证 $R^n \subseteq R$，从而
      $$R^m \circ R^n \subseteq R \circ R \subseteq R$$
      故 $R^{m+n} \subseteq R$，即 $R^n$ 传递。
    \end{subanswer}
  \end{subquestion}
\end{question}

\end{document}
